\chapter{项目五：异步日志器}

\section{项目目标}

日志器是一个综合项目：它接收多线程提交的日志消息，后台线程写入文件，关闭时刷新剩余消息。这个项目覆盖字符串格式化、时间、互斥量、条件变量、文件 RAII、异常边界和背压策略。

目标接口：

\begin{lstlisting}[language=C++,caption={日志器接口}]
class AsyncLogger {
public:
    explicit AsyncLogger(std::filesystem::path path);
    ~AsyncLogger();

    void info(std::string message);
    void error(std::string message);
};
\end{lstlisting}

接口按值接收 \code{std::string}，因为日志器要保存消息到后台队列。调用者传临时字符串时可以移动，传左值时复制一份，语义清楚。

\section{日志消息类型}

\begin{lstlisting}[language=C++,caption={日志消息}]
enum class LogLevel {
    info,
    error,
};

struct LogMessage {
    LogLevel level = LogLevel::info;
    std::chrono::system_clock::time_point time;
    std::string text;
};
\end{lstlisting}

不要把日志级别保存成裸字符串。枚举让内部逻辑更稳定，输出时再转换成文本。

\section{从同步日志推到异步日志}

先写最直接的同步版本：

\begin{lstlisting}[language=C++,caption={同步日志器}]
class SyncLogger {
public:
    explicit SyncLogger(std::filesystem::path path)
        : output_{path}
    {
        if (!this->output_) {
            throw std::runtime_error{"cannot open log file"};
        }
    }

    void info(std::string_view message)
    {
        std::lock_guard<std::mutex> lock{this->mutex_};
        this->output_ << "INFO " << message << '\n';
    }

private:
    std::ofstream output_;
    std::mutex mutex_;
};
\end{lstlisting}

这个版本容易验证：调用 \code{info} 时，当前线程直接写文件。缺点也明显：文件 I/O 慢时，业务线程会卡在 \code{info} 里。异步日志器不是为了让代码更花，而是把慢 I/O 从调用线程移到后台线程。

从同步到异步，本质上多了三件事：

\begin{enumerate}
  \item 调用线程只构造消息并放入队列。
  \item 后台线程从队列取消息并写文件。
  \item 关闭时要让后台线程写完剩余消息再退出。
\end{enumerate}

所以异步日志器首先是一个队列协议项目，其次才是格式化项目。

\section{队列和关闭协议}

日志器内部和任务队列类似：

\begin{lstlisting}[language=C++,caption={日志器内部状态}]
class AsyncLogger {
public:
    explicit AsyncLogger(std::filesystem::path path);
    ~AsyncLogger();

    void info(std::string message);
    void error(std::string message);

private:
    void push(LogLevel level, std::string message);
    void run();

    std::ofstream output_;
    std::mutex mutex_;
    std::condition_variable ready_;
    std::queue<LogMessage> messages_;
    bool closed_ = false;
    std::jthread worker_;
};
\end{lstlisting}

成员顺序要注意：\code{worker_} 最后声明，析构时最先析构并 join。但析构函数会先设置关闭状态并通知线程，让后台线程自然退出。

\section{关闭协议的状态推导}

关闭不是简单把 \code{closed_} 设成 \code{true}。要先定义规则：

\begin{itemize}
  \item 关闭前已经进入队列的消息必须写完。
  \item 关闭后新的 \code{info} 和 \code{error} 不再进入队列。
  \item 后台线程在队列为空且已关闭时退出。
  \item 析构函数不能抛异常。
\end{itemize}

由这四条规则推出后台线程的退出条件：

\begin{lstlisting}[language=C++,caption={退出条件由关闭规则推出}]
this->ready_.wait(lock, [this] {
    return this->closed_ || !this->messages_.empty();
});
if (this->messages_.empty() && this->closed_) {
    return;
}
\end{lstlisting}

如果只写 \code{if (this->closed_) return;}，会丢掉关闭前已经排队的消息。正确判断是“队列空并且关闭”才退出。这个边界是异步日志器最容易写错的地方。

\section{构造和析构}

\begin{lstlisting}[language=C++,caption={构造和析构日志器}]
AsyncLogger::AsyncLogger(std::filesystem::path path)
    : output_{path}, worker_{[this] { this->run(); }}
{
    if (!this->output_) {
        throw std::runtime_error{"cannot open log file"};
    }
}

AsyncLogger::~AsyncLogger()
{
    {
        std::lock_guard<std::mutex> lock{this->mutex_};
        this->closed_ = true;
    }
    this->ready_.notify_all();
}
\end{lstlisting}

这里有一个细节：构造函数先初始化 \code{output_}，再启动 \code{worker_}。成员按声明顺序初始化，而不是按初始化列表顺序。构造期间启动线程要非常小心，确保线程访问的成员已经构造完成。

\section{构造函数里启动线程的风险}

更稳的构造方式是先打开文件，确认成功后再启动线程。因为 \code{std::jthread} 一旦构造，后台函数就可能立刻运行：

\begin{lstlisting}[language=C++,caption={显式 start 降低构造风险}]
class AsyncLogger {
public:
    explicit AsyncLogger(std::filesystem::path path)
        : output_{path}
    {
        if (!this->output_) {
            throw std::runtime_error{"cannot open log file"};
        }
        this->worker_ = std::jthread{[this] { this->run(); }};
    }

private:
    std::ofstream output_;
    std::mutex mutex_;
    std::condition_variable ready_;
    std::queue<LogMessage> messages_;
    bool closed_ = false;
    std::jthread worker_;
};
\end{lstlisting}

这种写法要求 \code{worker_} 可以默认构造，然后在构造函数体里赋值启动。好处是失败路径更清楚：文件打不开时对象构造失败，后台线程根本没有启动。教学项目可以先用这种形态，降低半构造对象被线程访问的风险。

\section{提交日志}

\begin{lstlisting}[language=C++,caption={提交日志消息}]
void AsyncLogger::info(std::string message)
{
    this->push(LogLevel::info, std::move(message));
}

void AsyncLogger::error(std::string message)
{
    this->push(LogLevel::error, std::move(message));
}

void AsyncLogger::push(LogLevel level, std::string message)
{
    {
        std::lock_guard<std::mutex> lock{this->mutex_};
        if (this->closed_) {
            return;
        }
        this->messages_.push(LogMessage{
            .level = level,
            .time = std::chrono::system_clock::now(),
            .text = std::move(message),
        });
    }
    this->ready_.notify_one();
}
\end{lstlisting}

关闭后提交日志，本项目选择静默忽略。也可以选择抛异常或返回 \code{bool}，但规则必须写进接口文档。析构期间再抛异常通常不是好主意。

\section{有界队列和丢弃计数}

无限队列把压力藏进内存。教学版可以增加容量，并定义简单策略：队列满时丢弃 \code{info}，尽量保留 \code{error}。为避免静默丢失，还要统计丢弃数量：

\begin{lstlisting}[language=C++,caption={带容量和丢弃计数的状态}]
class AsyncLogger {
private:
    static constexpr std::size_t LOGGER_MAX_QUEUE_SIZE = 8192;

    std::queue<LogMessage> messages_;
    std::uint64_t dropped_info_count_ = 0;
};
\end{lstlisting}

\code{push} 的策略可以写成明确分支：

\begin{lstlisting}[language=C++,caption={队列满时丢弃 info}]
void AsyncLogger::push(LogLevel level, std::string message)
{
    {
        std::lock_guard<std::mutex> lock{this->mutex_};
        if (this->closed_) {
            return;
        }
        if (this->messages_.size() >= LOGGER_MAX_QUEUE_SIZE &&
            level == LogLevel::info) {
            ++this->dropped_info_count_;
            return;
        }
        this->messages_.push(LogMessage{
            .level = level,
            .time = std::chrono::system_clock::now(),
            .text = std::move(message),
        });
    }
    this->ready_.notify_one();
}
\end{lstlisting}

这段代码还没有解决“队列满但来了 \code{error}”时怎么办。可以选择阻塞、同步写入、或者丢掉最旧的 \code{info} 给 \code{error} 腾位置。教材里故意把这个问题暴露出来：背压策略必须按业务定义，不能靠容器自动决定。

\section{后台写入}

\begin{lstlisting}[language=C++,caption={后台写日志}]
void AsyncLogger::run()
{
    while (true) {
        std::optional<LogMessage> message;
        {
            std::unique_lock<std::mutex> lock{this->mutex_};
            this->ready_.wait(lock, [this] {
                return this->closed_ || !this->messages_.empty();
            });
            if (this->messages_.empty()) {
                return;
            }
            message = std::move(this->messages_.front());
            this->messages_.pop();
        }
        this->write_message(*message);
    }
}
\end{lstlisting}

取出消息后立刻释放锁，再写文件。文件 I/O 可能慢，如果持锁写文件，提交日志的线程也会被阻塞。锁只保护队列，不保护所有操作。

\section{flush 的完整实现思路}

为了让测试和调用者知道“当前已经提交的日志写完了”，可以把队列项从单一 \code{LogMessage} 扩展成命令：

\begin{lstlisting}[language=C++,caption={日志队列命令}]
struct FlushCommand {
    std::shared_ptr<std::promise<void>> done;
};

using LoggerCommand = std::variant<LogMessage, FlushCommand>;
\end{lstlisting}

队列从 \code{std::queue<LogMessage>} 改成 \code{std::queue<LoggerCommand>}。提交 \code{flush}：

\begin{lstlisting}[language=C++,caption={提交 flush 命令}]
void AsyncLogger::flush()
{
    auto done = std::make_shared<std::promise<void>>();
    std::future<void> future = done->get_future();
    {
        std::lock_guard<std::mutex> lock{this->mutex_};
        if (this->closed_) {
            return;
        }
        this->messages_.push(FlushCommand{.done = done});
    }
    this->ready_.notify_one();
    future.get();
}
\end{lstlisting}

后台线程处理命令：

\begin{lstlisting}[language=C++,caption={后台处理 flush 命令}]
void AsyncLogger::handle_command(LoggerCommand command)
{
    std::visit([this](auto& item) {
        using Item = std::decay_t<decltype(item)>;
        if constexpr (std::is_same_v<Item, LogMessage>) {
            this->write_message(item);
        } else {
            this->output_.flush();
            item.done->set_value();
        }
    }, command);
}
\end{lstlisting}

这里的顺序很关键：\code{flush} 命令在队列中的位置保证它前面的日志消息已经被处理；\code{output_.flush()} 保证文件缓冲被刷新；\code{promise} 保证调用线程有明确等待点。测试就可以调用 \code{flush()}，再读取文件断言内容，而不是睡一段时间碰运气。

\section{格式化输出}

\begin{lstlisting}[language=C++,caption={格式化日志级别}]
std::string_view level_name(LogLevel level)
{
    switch (level) {
    case LogLevel::info:
        return "INFO";
    case LogLevel::error:
        return "ERROR";
    }
    return "UNKNOWN";
}
\end{lstlisting}

写消息：

\begin{lstlisting}[language=C++,caption={写日志消息}]
void AsyncLogger::write_message(const LogMessage& message)
{
    this->output_ << level_name(message.level)
                  << " "
                  << message.text
                  << '\n';
}
\end{lstlisting}

时间格式化可以继续扩展。第一版先保证线程和关闭协议正确，不要把所有功能一次塞进去。

\section{背压和丢日志策略}

当前日志器队列无限增长。如果日志产生速度超过写盘速度，会占用越来越多内存。可以增加最大队列长度：

\begin{lstlisting}[language=C++,caption={日志队列背压策略}]
constexpr std::size_t LOGGER_MAX_QUEUE_SIZE = 8192;
\end{lstlisting}

队列满时常见策略：

\begin{itemize}
  \item 阻塞调用者：日志不能丢，但业务线程可能变慢。
  \item 丢弃新日志：业务不阻塞，但可能缺关键信息。
  \item 丢弃低级别日志：保留错误，丢掉 info。
  \item 同步写入：队列满时调用者自己写文件，复杂度更高。
\end{itemize}

日志系统没有完美策略，只有适合业务的取舍。教学版可以先选“队列满时丢弃 info，保留 error”，并统计丢弃数量。

\section{最终验收标准}

这个项目完成后，不应该只看“能打印几行日志”。要按协议验收：

\begin{enumerate}
  \item 构造失败时不启动后台线程。
  \item 多线程提交后，调用 \code{flush()} 能稳定看到已提交日志。
  \item 析构时队列中已有消息被写完。
  \item 关闭后提交不会抛出异常，也不会写入新消息。
  \item 队列满时的行为和文档一致，丢弃计数可观察。
  \item 后台写文件时不持有队列锁。
\end{enumerate}

这些标准都是从异步对象的状态机推出来的。项目代码再小，只要跨线程、跨生命周期，就要用协议而不是运气来设计。

\section{测试计划}

日志器测试不应该依赖睡眠。可以把后台写入逻辑拆成可测试的 \code{LogSink}，或者给日志器提供 \code{flush}：

\begin{itemize}
  \item 单线程写入一条 info，文件中出现一行。
  \item 多线程提交固定数量日志，最终行数正确。
  \item 析构后剩余消息被写完。
  \item 文件打不开时构造失败。
  \item 队列满时策略符合文档。
\end{itemize}

并发测试要避免“睡 100ms 应该写完”。更可靠的是用 \code{std::promise}、\code{std::latch} 或显式 \code{flush} 建立同步点。

\begin{exercisebox}
给异步日志器增加 \code{flush()}：调用后等待当前已经提交的日志全部写完。不要用睡眠。提示：可以向队列提交一个特殊消息，后台线程处理到它时通知 \code{promise}。
\end{exercisebox}
